% !TeX program = pdflatex
% papers/geometrico.tex — A RETA CONSTRUÍDA.
%
% Reescrito por inteiro (16/08). O paper antigo partia da dimensão 3, do círculo
% e do polígono, e construía uma «reta real geométrica» com a régua já na mão.
% O eval do coordenador inverteu a ordem:
%
%   «não usamos a reta para construir os números; construímos a reta a partir
%    do processo que produz os números»
%
% e fixou a espinha: bit → 8 leis → F_2^8 → dual → metálico/Pisot → corte → R,
% com a regra editorial de que NENHUMA definição geométrica precede o corte.
% Auto-contido: pdflatex geometrico.tex
\documentclass[11pt,a4paper]{article}
\usepackage[utf8]{inputenc}
\usepackage[T1]{fontenc}
\usepackage[portuguese]{babel}
\usepackage{amsmath,amssymb,amsthm}
\usepackage[margin=2.5cm]{geometry}
\usepackage{booktabs}
\usepackage{xcolor}
\usepackage[hidelinks]{hyperref}
\newtheorem{teorema}{Teorema}
\newtheorem{proposicao}[teorema]{Proposição}
\newtheorem{corolario}[teorema]{Corolário}
\newtheorem{definicao}{Definição}
\newtheorem{observacao}{Observação}
\newcommand{\code}[1]{\texttt{\small #1}}
\newcommand{\abs}[1]{\lvert #1\rvert}
\newcommand{\dg}{^{\dagger}}
% ─────────────────────────────────────────────────────────────────────────────
%  papers/gkcapa.tex — a capa do Reino, mínima, para os papers em `article`.
%
%  O estilo.tex completo é do `report` (títulos de capítulo, skak, …). Aqui fica
%  só o que a capa precisa, para o paper continuar a compilar sozinho com
%  pdflatex. O tradutor próprio (tests/tex) carrega o estilo.tex da raiz / o
%  symlink papers/estilo.tex e avalia o mesmo `\gkcapa`.
% ─────────────────────────────────────────────────────────────────────────────
\definecolor{ouro}{HTML}{B8912F}
\definecolor{ouroclaro}{HTML}{F3E9CE}
\definecolor{tinta}{HTML}{1E1B16}
\definecolor{regua}{HTML}{6E6858}
\providecommand{\gknota}{\fontsize{7.62}{11.05}\selectfont}
\providecommand{\gktexto}{\fontsize{10.50}{15.22}\selectfont}
\providecommand{\gksub}{\fontsize{12.33}{17.87}\selectfont}
\providecommand{\gksec}{\fontsize{14.47}{20.98}\selectfont}
\providecommand{\gkcap}{\fontsize{16.99}{24.63}\selectfont}
\providecommand{\gktit}{\fontsize{23.42}{33.95}\selectfont}
\providecommand{\Prj}{\mathbb{P}}
\providecommand{\Trf}{\mathcal{T}}
\providecommand{\gkcapa}[3]{%
\title{\vspace{-0.6cm}
{\color{ouro}\rule{\textwidth}{1.4pt}}\\[6mm]
\textbf{\gktit\color{tinta} Reino Dourado}\\[3mm]{\gksec\itshape\color{regua} Golden Kingdom}\\[8mm]
{\gkcap\scshape\color{ouro} #1}\\[5mm]
{\gksub\color{regua} #2}\\[2mm]
{\gktexto\color{regua}\itshape #3}\\[7mm]
{\color{ouro}\rule{\textwidth}{0.6pt}}\\[9mm]{\gknota\color{regua}\begin{minipage}{\textwidth}\setlength{\parindent}{0pt}%
\textbf{Aviso de Isenção Algorítmica:} O universo, as leis e as entidades contidas nestas páginas ---
incluindo felinos matriciais, aves de rapina algébricas e amuletos de controle de fluxo --- são
construções de pura ficção formal. Esta obra não descreve a realidade física, não serve como
engenharia reversa do cosmos e não constitui doutrina religiosa. Qualquer semelhança com bugs reais do seu
sistema operacional é mera coincidência (ou falta de reza). Prossiga por sua própria conta e
risco.\end{minipage}}}
\author{}\date{}}

\begin{document}
\gkcapa{A Reta Construída}
       {do bit às oito leis, do dual à família metálica, de Pisot ao corte}
       {a completude não é axioma: é a área que se conserva}
\maketitle

\begin{abstract}
\noindent Este paper não usa a reta para construir os números. Constrói a reta
a partir do processo que os produz, e a ordem é
\[
\boxed{\;
\text{bit}\;\to\;8\ \text{leis}\;\to\;\mathbb{F}_2^{8}\;\to\;\text{dual}
\;\to\;\text{metálico/Pisot}\;\to\;\text{corte}\;\to\;\mathbb{R}.\;}
\]
A base das oito leis é ortonormal, e por isso \emph{ler uma coordenada é medir
um bit}: a geometria discreta existe antes de qualquer real. O dual
$\nu(x)=-1/x$ produz a recorrência inteira $u_{n+2}=m\,u_{n+1}+u_n$, e com ela a
família metálica --- \emph{sem avaliar uma raiz}. Como
$\abs{\sigma_m}>1$ e $\abs{\sigma_m\dg}<1$, e como
$\sigma_m^{k}+(\sigma_m\dg)^{k}$ é \textbf{inteiro},
\[
\boxed{\;\operatorname{dist}(\sigma_m^{k},\mathbb{Z})=\abs{\sigma_m\dg}^{k}
\;:\;\text{a aproximação não é escolhida --- é produzida pela dualidade.}\;}
\]
Daí saem os intervalos encaixados e o corte. E o corte fecha por uma razão que
já é desta casa: $\abs{\det}=\abs{\sigma\sigma\dg}=1$, a área conserva-se; como
o comprimento diverge, a largura \emph{tem} de ir a zero. Só depois disso
entram comprimento, distância, reta e círculo --- e $\pi$ fica no lugar certo:
$\pi_k$ é a representação fechada do andar~$k$, e $\pi_\infty$ é o limite
parabólico, que não fecha na escada. \emph{O círculo não é o ponto de partida
da geometria: é o gume da construção.}

\smallskip
\noindent\emph{Medido:} \code{tests/geometria\_real.c} ($9{:}0$), com
\code{tests/metalica.c}, \code{tests/ortonormal.c}, \code{tests/pisotbase.c},
\code{tests/ordenado.c}, \code{tests/pi\_familia.js},
\code{tests/arquimedes\_area.js} e \code{tests/geometrico.js}.
\end{abstract}

%======================================================================
\section{A regra editorial, e o que ela proíbe}

Um paper que constrói $\mathbb{R}$ e usa a palavra \emph{comprimento} na
página um já pediu emprestado o que ia construir. Por isso a regra, e ela
governa a ordem das secções:

\begin{center}
\fbox{\parbox{0.9\textwidth}{\centering
\textbf{Nenhuma definição geométrica antes da construção do corte}
--- salvo as mínimas para enunciar o que será construído.}}
\end{center}

Concretamente: até ao \S\ref{sec:corte} não há reta, não há círculo, não há
polígono. Há bits, há uma forma bilinear sobre $\mathbb{F}_2$, há uma
recorrência de inteiros e há comparações de inteiros. As
secções~\ref{sec:geo} em diante é que trazem a geometria --- e trazem-na
\emph{como consequência}, não como cenário.

\medskip
\noindent\textbf{E a regra admite exactamente duas excepções, que se nomeiam
aqui para que não passem caladas.} Uma regra que o próprio texto viola sem
avisar é pior do que nenhuma.
\begin{itemize}
\item \textbf{A distância no cubo} (\S\ref{sec:orto}) --- e ela não é um
empréstimo, porque toma valores em $\{0,\ldots,8\}$: é uma contagem de bits.
Que exista \emph{geometria discreta} antes de qualquer real é o primeiro
resultado do paper, não uma antecipação do último.
\item \textbf{A área, o comprimento e a largura} do \S\ref{sec:area} --- e aí
as três palavras nomeiam \emph{um determinante}, que numa matriz de inteiros é
um inteiro. É a leitura geométrica de $\det$, e é o Teorema do Gato: nenhum
contínuo entra para a dizer. A distinção está enunciada em
Obs.~\ref{obs:area-inteira}, e é ela que impede a demonstração de ser
circular.
\end{itemize}

\medskip
\noindent Isto também arruma a divisão de trabalho com o resto do quadro:
\[
\boxed{\;
\underbrace{\text{\emph{Corpo Universal}}}_{\text{prova e organiza}}
\qquad\Big|\qquad
\underbrace{\text{este paper}}_{\text{realiza geometricamente a MESMA construção}}
\;}
\]
Nada aqui é um segundo sistema. O que aqui se acrescenta é a leitura
geométrica de uma cadeia que já está provada, e um mecanismo --- a conservação
da área --- que explica \emph{por que} o último elo fecha.

%======================================================================
\section{A primitiva: um bit}\label{sec:bit}

\begin{definicao}[o objecto mínimo]\label{def:bit}
$\mathbb{F}_2=\{0,1\}$, com a soma XOR e o produto AND. O oposto de $x$ é o
próprio $x$: \emph{o sinal não existe}.
\end{definicao}

\begin{definicao}[as oito posições]\label{def:posicoes}
$b=(b_0,\ldots,b_7)\in\mathbb{F}_2^{8}$, e a posição $i$ está
\textbf{reservada à Lei}~$i$. Escrito em binário, $b=\sum_i b_i 2^{i}$: um byte.
\end{definicao}

Que as leis sejam oito e que o byte tenha oito posições não é coincidência de
máquina: a dobra leva $1\to2\to4\to8$ e o catálogo das leis é um ciclo gerador
de período oito. Isso está estabelecido no \emph{Corpo Universal} e aqui
\emph{usa-se}, não se repete.

%======================================================================
\section{A base é ortonormal --- e por isso já há geometria}\label{sec:orto}

Sobre $\mathbb{F}_2$ há uma só forma bilinear simétrica não degenerada
disponível, e é esta:
\[
\langle a,b\rangle \;=\; \operatorname{paridade}(a\wedge b).
\]

\begin{teorema}[a base das oito leis é ortonormal]\label{thm:orto}
$\langle e_i,e_j\rangle=\delta_{ij}$, isto é, a matriz de Gram é a identidade.
Logo a coordenada recupera-se sem inverter sistema nenhum:
\[
\boxed{\;b_i=\langle b,e_i\rangle\;}
\]
--- \emph{um AND e uma paridade, que é ler o bit}. \emph{Medido:}
\code{tests/ortonormal.c} ($6{:}0$). $\square$
\end{teorema}

\noindent E é aqui que a primeira frase geométrica aparece, ainda sem um único
real dentro:

\begin{proposicao}[a distância é a soma das coordenadas]\label{prop:dist}
Em $\mathbb{F}_2^{8}$,
\[
d(a,b)\;=\;\sum_{i=0}^{7}\bigl\lvert\langle a,e_i\rangle-\langle b,e_i\rangle\bigr\rvert
\;=\;\operatorname{peso}(a\oplus b),
\]
e $d$ toma valores em $\{0,\ldots,8\}$, anula-se só na diagonal, é simétrica e
fecha a desigualdade do triângulo. O \textbf{diâmetro é $8$}, porque as leis são
oito. \emph{Medido} nos $65\,536$ pares e nos $262\,144$ trios, pelos dois
caminhos --- a soma das coordenadas e o peso ---, sem divergência:
\code{tests/geometria\_real.c} \S G0. $\square$
\end{proposicao}

\begin{observacao}[o que já se ganhou, e o que ainda falta]
Existe distância, existe diâmetro, existe triângulo. Ou seja: \emph{há
geometria}, e ela é inteira, finita e exacta. O que ainda \textbf{não} há é
\emph{contínuo} --- e é exactamente isso que as secções seguintes constroem.
Esta face, aliás, não admite ordem de corpo, porque $1+1=0$; dizer o contrário
seria trocar o escopo, e o escopo faz parte do enunciado.
\end{observacao}

%======================================================================
\section{O dual: uma unidade, duas faces}\label{sec:dual}

A dualidade organiza as oito posições sem acrescentar graus de liberdade. É o
princípio que no quadro aparece como Gato e Esquilo:
\[
\boxed{\;\text{uma unidade}\;=\;\text{duas faces duais.}\;}
\]
Não é «bit \emph{versus} real»: é a mesma estrutura vista dos dois lados. E o
que interessa aqui é que o dual não é decorativo --- ele \textbf{produz} o
objecto seguinte.

\begin{definicao}[a involução]\label{def:nu}
$\nu(x)=-1/x$. Ela é total na recta projectiva, onde a inversão é a troca
$[p:q]\mapsto[q:p]$ e $0\dg=\infty$.
\end{definicao}

\begin{teorema}[o dual produz a recorrência]\label{thm:dual-produz}
Seja $A_m=\begin{pmatrix}m&1\\1&0\end{pmatrix}$ e $\sigma$ a classe de $x$ em
$\mathbb{Z}[x]/(x^{2}-mx-1)$. Então $\sigma\dg=-\sigma^{-1}=m-\sigma$, e
\[
\sigma+\sigma\dg=m,\qquad \sigma\,\sigma\dg=-1,
\]
lidos por \textbf{três caminhos independentes} que têm de concordar:
\emph{(A)} o traço e o determinante de $A_m$;
\emph{(B)} a involução $\nu$ dentro de $\mathbb{Z}[\sigma]$, usando apenas
$\sigma^{2}=m\sigma+1$;
\emph{(C)} a potência
$A_m^{k}=\begin{pmatrix}F_{k+1}&F_k\\F_k&F_{k-1}\end{pmatrix}$.
Concordam para $m=1,\ldots,40$; e o \textbf{gume}: alterada uma entrada da
companheira, o determinante deixa de valer $-1$ nos $40$ casos.
\emph{Medido:} \code{tests/geometria\_real.c} \S G1. $\square$
\end{teorema}

\emph{Nenhuma raiz foi avaliada.} As duas identidades acima são a soma e o
produto de duas faces, e saem da redução --- não de um valor decimal.

%======================================================================
\section{A escada metálica: a recorrência é a construção}\label{sec:metal}

\begin{definicao}[a escada]\label{def:escada}
Para $m\geq1$ inteiro,
\[
u_{n+2}=m\,u_{n+1}+u_n,
\]
e escrevem-se $F_k$ (com $F_0=0$, $F_1=1$) e $t_k$ (com $t_0=2$, $t_1=m$) as
duas soluções que a casa usa: o \emph{Fibonacci} e o \emph{Lucas} metálicos.
\end{definicao}

A representação alcançada é $\sigma_m=\bigl(m+\sqrt{m^{2}+4}\bigr)/2$, e ela
\textbf{não} é o ponto de partida: escreve-se aqui uma vez, para dizer de que se
fala, e não volta a ser usada como número. \emph{A raiz é a representação
alcançada; a recorrência é a construção.} Tudo o que segue lê-se em
$\mathbb{Z}[\sigma]$.

\begin{proposicao}[o que já estava medido]\label{prop:ja}
$\sigma^{k}=F_{k-1}+F_k\sigma$; o traço $t_k=\sigma^{k}+(\sigma\dg)^{k}$ fica em
$\mathbb{Z}$ e obedece a $t_{k+1}=m\,t_k+t_{k-1}$; $N(\sigma^{k})=(-1)^{k}$; e
\emph{unidade} e \emph{Pisot} são critérios distintos, com a família metálica na
intersecção. \emph{Medido:} \code{tests/metalica.c}. $\square$
\end{proposicao}

%======================================================================
\section{Pisot é o mecanismo geométrico}\label{sec:pisot}

Este é o coração do paper.

\begin{teorema}[a aproximação é produzida, não escolhida]\label{thm:pisot-geo}
Como $\abs{\sigma_m}>1$ e $\abs{\sigma_m\dg}<1$, tem-se
$(\sigma_m\dg)^{k}\to0$; e como $\sigma_m^{k}+(\sigma_m\dg)^{k}=t_k$ é
\textbf{inteiro},
\[
\boxed{\;\operatorname{dist}\bigl(\sigma_m^{k},\mathbb{Z}\bigr)
=\bigl\lvert(\sigma_m\dg)^{k}\bigr\rvert\;}
\]
sempre que $\abs{\sigma_m\dg}^{k}<\tfrac12$ --- condição que se decide em
inteiros puros, por
\[
(t_k-1)^{2}\;<\;\Delta F_k^{2}\;<\;(t_k+1)^{2},\qquad \Delta=m^{2}+4,
\]
e que, substituída a identidade $t_k^{2}-\Delta F_k^{2}=4(-1)^{k}$, se reduz a
$t_k\geq2$ ($k$ ímpar) e $t_k\geq3$ ($k$ par). \emph{Os dois caminhos concordam
em todos os $131$ andares medidos, sem uma discordância.}
\emph{Medido:} \code{tests/geometria\_real.c} \S G2; e
\code{tests/pisotbase.c} \S B4 pelo lado da numeração. $\square$
\end{teorema}

\begin{observacao}[a hipótese não é vazia --- e o escopo diz-se]\label{obs:gume-ouro}
Há \textbf{exactamente um} andar, em toda a janela medida, onde
$\abs{\sigma\dg}^{k}<\tfrac12$ falha: o \emph{ouro no primeiro degrau}. Lá,
$t_1=1$, $\Delta F_1^{2}=5$ e $(t_1+1)^{2}=4$ --- e de facto o inteiro mais
próximo de $\varphi$ é $2$, não $1$. Portanto a lei vale do segundo degrau em
diante para $m=1$, e desde o primeiro para $m\geq2$. Uma afirmação sem este
escopo estaria meio certa, que é a pior maneira de estar certo.
\end{observacao}

\medskip
\noindent A frase que este teorema autoriza é a do coordenador, e ela é a tese
do paper:
\[
\boxed{\;\text{a aproximação não é escolhida; ela é PRODUZIDA pela dualidade.}\;}
\]
Nenhuma sucessão foi exibida à mão, nenhum $\varepsilon$ foi escolhido. O erro
que aparece é $\abs{\sigma\dg}^{k}$, e $\sigma\dg$ é a \emph{outra face} de
$\sigma$ --- a que o \S\ref{sec:dual} produziu.

%======================================================================
\section{O encaixe}\label{sec:encaixe}

\begin{teorema}[os convergentes encaixam, e a largura é inteira]\label{thm:encaixe-geo}
Seja $c_k=F_{k+1}/F_k$. Então:
\emph{(i)} $c_k$ e $c_{k+1}$ caem em lados \textbf{opostos} de $\sigma$;
\emph{(ii)} cada intervalo está \textbf{dentro} do anterior;
\emph{(iii)} pela identidade de Cassini metálica
$F_{k+1}^{2}-F_{k+2}F_k=(-1)^{k}$, a largura é
\[
\abs{c_{k+1}-c_k}=\frac{1}{F_kF_{k+1}}
\]
--- exacta, e obtida sem fazer uma única divisão. \emph{Medido} nos $115$
degraus de $m=1,\ldots,8$, sem excepção: \code{tests/geometria\_real.c} \S G3.
$\square$
\end{teorema}

\begin{observacao}
A primeira versão deste enunciado tinha o sinal de Cassini trocado, e a medição
deu $0$ de $115$. Um sinal errado num expoente $(-1)^{k}$ é invisível na
leitura e total na medição --- é para isso que o medidor serve.
\end{observacao}

%======================================================================
\section{A área é a razão: por que o encaixe fecha}\label{sec:area}

Aqui entra o Teorema do Gato, e é ele que substitui o axioma.

\begin{teorema}[a completude é consequência da conservação da área]\label{thm:area-fecha}
$\det A_m=-1$, logo $\det A_m^{k}=(-1)^{k}$ e
\[
\boxed{\;\abs{\det}=\abs{\sigma\sigma\dg}=1\quad\text{em todos os andares.}\;}
\]
Lido geometricamente: o quadrado unitário vai, sob $A_m^{k}$, num paralelogramo
de \textbf{comprimento} $\sigma^{k}$ na direcção que estica e \textbf{largura}
$\abs{\sigma\dg}^{k}$ na que encolhe, e
\[
\text{comprimento}\times\text{largura}=1 .
\]
Como o comprimento diverge, a largura \textbf{tem} de ir a zero. \emph{Não é um
axioma acrescentado à construção: é a área que não pode aumentar nem diminuir.}
Medido: $\abs{\det}=1$ nos $115$ andares, e o traço a crescer nos $115$.
\emph{Medido:} \code{tests/geometria\_real.c} \S G4. $\square$
\end{teorema}

\begin{observacao}[a «área» aqui é um inteiro --- e é por isso que o argumento
não é circular]\label{obs:area-inteira}
Um paper que constrói $\mathbb{R}$ e usa \emph{área} para o fazer devia estar a
pedir emprestado o que vai construir. Não está, e a razão é aritmética: a
quantidade conservada é
\[
\det A_m^{k}=F_{k+1}F_{k-1}-F_k^{2}=(-1)^{k},
\]
uma \textbf{diferença de produtos de inteiros}. «Área» é o \emph{nome
geométrico} desse inteiro --- o Teorema do Gato: $\det$ \textbf{é} a medida ---,
e a igualdade
$\text{comprimento}\times\text{largura}=1$ é a factorização de $\pm1$ nas duas
faces $\sigma$ e $\sigma\dg$. Nenhum contínuo foi usado para a enunciar nem para
a medir: o medidor compara inteiros. O contínuo aparece \emph{depois}, e
aparece como \textbf{conclusão} --- a largura vai a zero --- e não como
hipótese.
\end{observacao}

\begin{observacao}[o único degrau plano]
$F$ cresce estritamente em $114$ dos $115$: a excepção é $F_1=F_2=1$, o ouro,
que demora um andar a arrancar. É o \emph{mesmo} $m$ do
Obs.~\ref{obs:gume-ouro} --- o primeiro degrau do ouro é o sítio onde a família
ainda não separou as duas faces o suficiente.
\end{observacao}

\begin{teorema}[o controlo: sem Pisot não há reta]\label{thm:controlo}
Tome-se $x^{2}-2x-4$, cuja raiz dominante é $1+\sqrt5$ e cuja conjugada é
$1-\sqrt5$, com $\abs{\sigma\dg}=\sqrt5>1$ e norma $-4$. A recorrência é
inteira, o traço é inteiro, e ainda assim:
o critério do Teor.~\ref{thm:pisot-geo} falha em \textbf{todos} os $12$ andares
medidos, e $\abs{\det}=1$ em \textbf{nenhum} --- $\det A^{k}=(-4)^{k}$. O que
devia encolher \emph{cresce}, e o encaixe não fecha.
\emph{Uma recorrência inteira qualquer não dá a reta: tem de ser unidade e tem
de ser Pisot.} \emph{Medido:} \code{tests/geometria\_real.c} \S G4. $\square$
\end{teorema}

\noindent Sem este controlo, o Teor.~\ref{thm:pisot-geo} poderia estar a medir
uma propriedade de \emph{qualquer} recorrência inteira. Está a medir Pisot.

%======================================================================
\section{O corte}\label{sec:corte}

\begin{teorema}[o corte decide-se em inteiros]\label{thm:corte}
Para $a\in\mathbb{Z}$, $b\geq1$ e $\Delta=m^{2}+4$:
\[
\frac{a}{b}<\sigma_m
\quad\Longleftrightarrow\quad
2a-mb<0 \ \ \text{ou}\ \ (2a-mb)^{2}<b^{2}\Delta .
\]
Nos $19\,360$ racionais varridos ($m\leq8$, $b\leq20$, $\abs{a}\leq60$):
\emph{(i)} cada um cai de um lado e \textbf{nenhum} cai em cima --- porque
$\Delta$ não é quadrado, que é dizer $\sigma\notin\mathbb{Q}$;
\emph{(ii)} nenhuma das duas classes tem extremo: para cada $a/b$ do lado de
baixo \emph{exibe-se} $a'/b'$ com $a/b<a'/b'<\sigma$, e simetricamente do outro.
A raiz nunca é avaliada: elevar ao quadrado é tudo o que é preciso.
\emph{Medido:} \code{tests/geometria\_real.c} \S G5. $\square$
\end{teorema}

\begin{teorema}[e o bit lê o corte]\label{thm:bit-le}
Bissectar $[m,m+1]$ e perguntar de que lado está $\sigma_m$ é \emph{uma}
pergunta --- a do Teor.~\ref{thm:corte}, decidida em inteiros --- e devolve
exactamente \textbf{um bit}. Oito passos enchem um byte; o intervalo fica com
largura $1/256$ e $\sigma_m$ continua dentro nos oito passos, para os oito
metais. E o bit lê-se pela coordenada $\langle b,e_i\rangle$ da base ortonormal
do \S\ref{sec:orto}. \emph{Medido:} \code{tests/geometria\_real.c} \S G6.
$\square$
\end{teorema}

\medskip
\noindent A cadeia fecha:
\[
\boxed{\;
\text{Pisot}\Longrightarrow\text{vazamento}\to0
\Longrightarrow\text{encaixe}\Longrightarrow\text{corte}\Longrightarrow\mathbb{R}.\;}
\]

\begin{observacao}[a identificação final é literatura, e diz-se]\label{obs:citada}
Que um corpo ordenado completo seja isomorfo a $\mathbb{R}$, e único a menos de
isomorfismo, é \textbf{teorema clássico} --- não deste quadro. O que aqui se
constrói e se mede são as \emph{hipóteses}: corpo, ordem compatível com as
operações, arquimediano com o $n$ exibido, e completo pelo corte. A conclusão
cita-se. Misturar as duas coisas seria dar ao trabalho um crédito que não é
dele. \emph{Medido:} \code{tests/ordenado.c} ($7{:}0$).
\end{observacao}

%======================================================================
\section{E só agora: a geometria}\label{sec:geo}

Com o corte construído, as definições geométricas podem entrar --- e entram
\emph{como consequências}, cada uma apoiada numa peça já medida.

\begin{definicao}[comprimento, intervalo, distância, reta]\label{def:geo}
O \textbf{comprimento} do intervalo $[c_k,c_{k+1}]$ é $1/(F_kF_{k+1})$
(Teor.~\ref{thm:encaixe-geo}); a \textbf{distância} entre dois cortes é o
comprimento do intervalo que os separa; e a \textbf{reta} é o conjunto dos
cortes, com a ordem induzida pela comparação inteira do Teor.~\ref{thm:corte}.
\end{definicao}

\noindent Repare-se no que foi invertido. No paper anterior, a reta era o
cenário e os números eram pontos marcados nela. Aqui a reta é o
\emph{arquivo do processo}: cada ponto é o corte que um encaixe produziu, e o
comprimento é a largura que a conservação da área obrigou a ir a zero.

%======================================================================
\section{O círculo é o gume}\label{sec:circulo}

Só agora o círculo, e ele entra pelo lado que interessa: como o \emph{membro que
não fecha}.

\begin{teorema}[$\pi_k$ fecha por andar; $\pi_\infty$ não fecha na escada]\label{thm:pi}
Pela tricotomia do traço, a família é uma só: os \textbf{polígonos} são os
membros \emph{elípticos} ($t_n<2$), e fecham em ordem $2n$ exacta no primo da
ordem --- $\sqrt2$ é $11$ em $\mathbb{F}_{17}$, $\sqrt3$ é $9$ em
$\mathbb{F}_{13}$; os \textbf{metais} são os \emph{hiperbólicos}
($\Delta=m^{2}+4>0$), que são os que este paper usou para construir o corte; e o
\textbf{círculo} é o limite \emph{parabólico} $t=2$, unipotente, cuja ordem
módulo $p$ é exactamente $p$ --- logo \textbf{nunca} um divisor da ordem da
escada. \emph{Medido:} \code{tests/pi\_familia.js} ($7{:}0$). $\square$
\end{teorema}

\begin{corolario}[o lugar de $\pi$]\label{cor:pi}
$\pi_k$ é a representação geométrica \textbf{fechada} do andar~$k$, e
$\pi_\infty$ é o limite parabólico, que não fecha em ordem da escada. Portanto
não há aqui nada a aproximar: quem escrever um limiar para $\pi$ está a importar
uma constante que o quadro já tem exacta por andar, ou está a pedir o limite. E
esta é a mesma fronteira do eixo de Pontryagin, agora dita como propriedade do
membro-limite: \emph{a álgebra opera e não alcança o contínuo; o encaixe alcança
e não opera}. $\square$
\end{corolario}

\begin{observacao}[o círculo como gume, e não como origem]
O paper anterior começava no triângulo --- «o primeiro polígono regular que
fecha» --- e subia. A ordem estava invertida: o polígono e o círculo são o
\emph{fim} da construção, não o princípio, porque ambos precisam de área, e a
área precisa de comprimento, e o comprimento precisa do corte. Aqui eles são o
gume: o sítio onde se vê que a escada tem uma fronteira, e que a fronteira é
teorema.
\end{observacao}

%======================================================================
\section{As áreas dos andares, e a torre normalizada}\label{sec:areas}

O que o paper anterior tinha de próprio sobrevive --- \emph{a jusante} do corte,
que é onde sempre pertenceu. Fica registado por inteiro, com os seus medidores.

\begin{proposicao}[Arquimedes das áreas]\label{prop:arquimedes}
Cada andar elíptico $n\in\{3,4,6,8\}$ traz um círculo renormalizado com a área
do polígono circunscrito $C_n$, ficando o inscrito $I_n$ por baixo; e a
duplicação $n\to2n$ desce pelas médias \emph{geométrica} e \emph{harmónica}:
\[
I_{2n}^{2}=I_nC_n,\qquad C_{2n}(I_{2n}+C_n)=2I_{2n}C_n .
\]
Os círculos dos andares \textbf{não são o mesmo círculo}: $C_3,C_4,C_6,C_8$ são
distintos dois a dois, por desigualdades quadráticas inteiras. E
$I_8<333/106<\pi<22/7<C_8$ em $\mathbb{Z}$: $\pi$ no buraco. \emph{O gume}: a
média \emph{aritmética} não duplica. \emph{Medido:}
\code{tests/arquimedes\_area.js} ($21{:}0$). $\square$
\end{proposicao}

\begin{proposicao}[a torre normalizada pesa $1$]\label{prop:torre}
Com $\abs{C_n}=\abs{P_n}=1/2^{\,n-1}$ para $n\geq3$, a soma parcial é
\[
\sum_{n=3}^{N}\frac{1}{2^{\,n-1}}=\frac12-\frac{1}{2^{\,N-1}},
\]
logo cada metade da torre pesa $\tfrac12$ e a torre pesa $1$; e a leitura em
andares consecutivos dá $\tfrac1{2^{n}}+\tfrac1{2^{n-1}}=\tfrac3{2^{n}}$ dos
dois lados. \emph{O gume}: repetido o peso do triângulo, três andares somam
$\tfrac34\neq\tfrac12$ --- a série só fecha porque os níveis \emph{descem} em
potência de $2$. \emph{Medido:} \code{tests/geometrico.js}. $\square$
\end{proposicao}

\begin{observacao}[o estatuto que estas duas mudaram]
Na versão anterior, a igualdade $\abs{C_n}=\abs{P_n}$ era o \emph{axioma
fundamental} e a normalização era a construção da reta. Nenhuma das duas coisas
se sustenta na ordem nova: a igualdade é uma \textbf{identificação de andar} e a
normalização é uma \textbf{escolha de pesos}, e ambas pressupõem a área, que
pressupõe o corte. Passam a proposições, e o que era o axioma passa a ser a
hipótese que as define. \emph{Nada foi apagado; foi mudado de andar} --- que é o
mesmo que se fez, no quadro, com a excepção do zero.
\end{observacao}

%======================================================================
\section{O que este paper não faz}\label{sec:escopo}

\begin{itemize}
\item \textbf{Não prova a unicidade de $\mathbb{R}$.} Cita-a
(Obs.~\ref{obs:citada}). O que mede são as hipóteses.
\item \textbf{Não constrói todos os reais}, e sim os de $\mathbb{Z}[\sigma_m]$:
a construção é dada na família metálica, que é onde o mecanismo --- Pisot --- é
exacto e inteiro. A extensão a um corte arbitrário é o Teor.~\ref{thm:corte},
que já decide qualquer racional, mas o \emph{gerador} aqui é metálico.
\item \textbf{Não afirma que a face finita é ordenada.} Não é
($1+1=0$), e isso é teorema. A face que alcança é que é o corpo ordenado
completo. O escopo é parte de cada enunciado.
\item \textbf{Não calcula $\pi$}, e não porque falte instrumento: porque
$\pi_k$ já é exacto por andar e $\pi_\infty$ é o limite, por
Teor.~\ref{thm:pi}.
\end{itemize}

\medskip
\begin{center}
\fbox{\parbox{0.92\textwidth}{\centering
\textbf{A tese.} A geometria real é a realização alcançada de uma estrutura
binária de oito leis: \textbf{a base fornece as coordenadas}, \textbf{a
dualidade fornece a recorrência}, \textbf{a família de Pisot fornece o encaixe}
e \textbf{o corte fornece o contínuo}.\\[4pt]
E o encaixe fecha porque a área se conserva: $\abs{\det}=\abs{\sigma\sigma\dg}=1$.}}
\end{center}

\medskip
\emph{Medido:} \code{tests/geometria\_real.c} ($9{:}0$),
\code{tests/ortonormal.c} ($6{:}0$), \code{tests/metalica.c},
\code{tests/pisotbase.c}, \code{tests/ordenado.c} ($7{:}0$),
\code{tests/pi\_familia.js} ($7{:}0$),
\code{tests/arquimedes\_area.js} ($21{:}0$), \code{tests/geometrico.js}.

\end{document}
